\scanpage{45}
\mspage{10}
2. \uline{Stratifications globales} : \uline{formalisme} $\frac12$ \uline{simplicial}
(sans tubes).

Pour simplifier, je vais me placer sur un
espace topologique $X$ --- par la suite $X$
deviendra un topos quelconque. Les construc-
tions qui suivent, \unsure{relatifs} à une ``stratifica-
tion globale'', devront ``faisceautiser sur $X$
par la suite, de la façon habituelle ---
ce qui amènera \inserted{alors} à imposer des condi-
tions supplémentaires de connexité et
de locale connexité, qui pour le moment
ne sont pas nécessaires. De même, il n'y
a pas lieu pour le moment de s'en-
combrer d'hypothèses de locale trivialité
normale des plongements d'espaces.

Soit $I$ un ensemble ordonné, \deleted{$\unsure{X_i}$}
\[
(1)\qquad (X_i)_{i\in I},\quad X_i\subset X
\]
une famille de \deleted{\unsure{parties}}\inserted{sous-espaces} de $X$. On suppose
\[
(1')\quad\left\{\begin{aligned}
\text{a)}\quad&\text{Les }X_i\text{ fermés, la famille }(X_i)\text{ localement finie.}\\
\text{b)}\quad&i\leq j\ \Longrightarrow\ X_i\subset X_j.
\end{aligned}\right.
\]
Posant
\[
(2)\qquad X_{\Delta_0}=\coprod_{i\in I}X_i
=\{(x,i)\in X\times I\mid x\in X_i\}
\]
on a un morphisme canonique
\mspage{44}
